\subsection{上同伦子}

设 H 为箭图，G 为群胚，\(\varphi\) 与 \(\psi\) 为从 H 到 G 的箭图态射。

以 \(G_{1}\) 表示按下述方式定义的箭图：\(G_{1}\) 的顶点就是 G 的顶点；\(G_{1}\) 的箭是集 \(\mathrm{Fl(G)}\) 与 \(\mathrm{Som(H)}\) 的不交并的元素；\(G_{1}\) 的源映射与 G 的源映射在 \(\mathrm{Fl(G)}\) 上相重合，且与 \(\mathrm{Som}(\varphi)\) 在 \(\mathrm{Som(H)}\) 上相重合；\(G_{1}\) 的靶映射与 G 的靶映射在 \(\mathrm{Fl(G)}\) 上相重合，且与 \(\mathrm{Som}(\psi)\) 在 \(\mathrm{Som(H)}\) 上相重合。以 \(\alpha_{1}\) 表示从 G 到 \(G_{1}\) 的箭图态射，它由 \(\mathrm{Som(G)}\) 的恒同映射以及 \(\mathrm{Fl(G)}\) 到 \(\mathrm{Fl(G_{1})}\) 中的典范单射定义。以 \(h_{1}\) 表示典范单射 \(\mathrm{Som(H)} \to \mathrm{Fl(G_{1})}\) 。

考虑自由群胚 \(\mathrm{Grp}(G_{1})\) （在 \(G_{1}\) 上构造；II, 第 174 页, 定义 9），并以 \(\theta_{1}\) 表示从 \(G_{1}\) 到 \(\mathrm{Grp}(G_{1})\) 的典范箭图态射。最后，以 \(\mathrm{Coh}(\varphi,\psi)\) 表示由 \(\mathrm{Grp}(G_{1})\) 经缩并圈（在 \(x\) 的起点处）而得的群胚。

\[
\alpha_ {1} (x) \alpha_ {1} (y) \alpha_ {1} (x y) ^ {- 1},\tag{1}
\]

对 G 的每一对可合成箭 \((x,y)\) ，以及圈（在 \(\varphi(a)\) 处）

\[
\alpha_ {1} (\varphi (f)) h _ {1} (b) \alpha_ {1} (\psi (f)) ^ {- 1} h _ {1} (a) ^ {- 1}\tag{2}
\]

其中 \(a\) 、\(b\) 属于 \(\operatorname{Som}(\mathrm{H})\) 且 \(f \in \operatorname{Fl}_{ab}(\mathrm{H})\) 。以 \(\pi: \operatorname{Grp}(\mathrm{G}_1) \to \operatorname{Coh}(\varphi, \psi)\) 表示典范态射；令 \(\alpha = \pi \circ \theta_1 \circ \alpha_1\) 与 \(h = \operatorname{Fl}(\pi \circ \theta_1) \circ h_1\) 。

命题 2. —— 群胚 Coh(φ,ψ) 由这样的子箭图生成：其顶点集为 Som(G)，其箭集为映射 Fl(α) 与 h 的像的并。

此子箭图是箭图 \(G_{1}\) 在箭图态射 \(\pi \circ \theta_{1}\) 下的像，故命题由 \(\mathrm{Coh}(\varphi, \psi)\) 的构造立即得出。

命题 3. —— 箭图态射 \(\alpha\) 是从 G 到 \(\mathrm{Coh}(\varphi, \psi)\) 的群胚态射，而映射 \(h\) 是连接 \(\alpha \circ \varphi\) 与 \(\alpha \circ \psi\) 的同伦。

三元组 \((\mathrm{Coh}(\varphi,\psi),\alpha,h)\) 具有下列泛性质：若 \(G'\) 为群胚，\(\alpha'\) 为从 G 到 \(G'\) 的群胚态射，且 \(h'\colon\mathrm{Som}(\mathrm{H})\to\mathrm{Fl}(\mathrm{G}')\) 为连接 \(\alpha'\circ\varphi\) 与 \(\alpha'\circ\psi\) 的同伦，则存在唯一的群胚态射 \(\eta\colon\mathrm{Coh}(\varphi,\psi)\to\mathrm{G}'\) ，使得

\[
\alpha^ {\prime} = \eta \circ \alpha \qquad\text{且}\qquad h ^ {\prime} = \operatorname{Fl} (\eta) \circ h.\tag{3}
\]

鉴于经缩并箭而得的群胚的定义，圈 (1) 的缩并蕴含 \(\alpha\) 是群胚态射，而圈 (2) 的缩并蕴含 \(h\) 是连接 \(\alpha \circ \varphi\) 与 \(\alpha \circ \psi\) 的同伦。

设 \(G'\) 、\(\alpha'\) 、\(h'\) 如命题所述。设 \(\eta_{1}\) 为从 \(G_{1}\) 到 \(G'\) 的箭图态射，使得 \(\operatorname{Som}(\eta_{1})\) 等于 \(\operatorname{Som}(\alpha')\) ，且 \(\operatorname{Fl}(\eta_{1})\) 与 \(\operatorname{Fl}(\alpha')\) 在 \(\operatorname{Fl}(G)\) 上相重合、与 \(h'\) 在 \(\operatorname{Som}(H)\) 上相重合。存在唯一的群胚态射 \(\eta_{2}: \operatorname{Grp}(G_{1}) \to G'\) ，使得 \(\eta_{1} = \eta_{2} \circ \theta_{1}\) 。由于 \(\alpha'\) 是群胚态射且 \(h'\) 是连接 \(\alpha' \circ \varphi\) 与 \(\alpha' \circ \psi\) 的同伦，\(\eta_{2}\) 通过取商定义了一个群胚态射 \(\eta\) （从 \(\operatorname{Coh}(\varphi, \psi)\) 到 \(G'\) ）（II, 第 170 页, 命题 3）。此态射满足关系式 (3)，且是唯一的（II, 第 185 页, 命题 2）。

定义 3. —— 群胚 Coh(φ,ψ) 称为偶 (φ,ψ) 的上同伦子。称 α 为从 G 到 Coh(φ,ψ) 的典范态射，称 h 为联系 α∘φ 与 α∘ψ 的典范同伦。

我们称偶 \((\varphi,\psi)\) 的框架为这样的箭图：其顶点集为 G 的轨道集，其箭集为 H 的连通分支集，其源映射与靶映射由 \(\varphi\) 与 \(\psi\) 通过取商诱导。

命题 4. —— 映射 Orb(α): Orb(G) → Orb(Coh(φ,ψ)) 是满射的，其纤维是偶 (φ,ψ) 的框架的连通分支。

态射 \(\alpha\) 是态射 \(\alpha_{1}\colon\mathrm{G}\to\mathrm{G}_{1}\) 、\(\theta_{1}\colon\mathrm{G}_{1}\to\mathrm{Grp}(\mathrm{G}_{1})\) 与 \(\pi\colon\mathrm{Grp}(\mathrm{G}_{1})\to\mathrm{Coh}(\varphi,\psi)\) 的合成。映射 \(\theta_{1}\) 诱导从箭图 \(\mathrm{G}_{1}\) 的连通分支集到 \(\mathrm{Grp}(\mathrm{G}_{1})\) 的轨道集的双射，且映射 \(\mathrm{Orb}(\pi)\colon\mathrm{Orb}(\mathrm{Grp}(\mathrm{G}_{1}))\to\mathrm{Orb}(\mathrm{Coh}(\varphi,\psi))\) 是双射（II, 第 170 页, 注 1）。因此只需证明从 \(\mathrm{Orb}(\mathrm{G})\) 到 \(\pi_{0}(\mathrm{G}_{1})\) 的映射（由 \(\alpha_{1}\) 诱导）是满射的，且其纤维是偶 \((\varphi,\psi)\) 的框架的连通分支。满射性由映射 \(\mathrm{Som}(\alpha_{1})\) 是恒同映射这一事实得出。\(\mathrm{Som}(\mathrm{G})\) 中由“\(a\) 与 \(b\) 位于 \(\mathrm{G}_{1}\) 的同一连通分支中”给出的等价关系，由关系“存在 G 的一条箭连接 \(a\) 与 \(b\)”以及“存在 H 的顶点 \(h\) 使得 \(\varphi(h)=a\) 且 \(\psi(h)=b\)”生成。此等价关系与从 \(\mathrm{Som}(\mathrm{G})\) 到 \(\mathrm{Orb}(\mathrm{G})\) 的映射相容，且在 \(\mathrm{Orb}(\mathrm{G})\) 上诱导的关系由关系“存在 H 的轨道 \(\eta\) 使得 \(\mathrm{Orb}(\varphi)(\eta)=\alpha\) 且 \(\mathrm{Orb}(\psi)(\eta)=\beta\)”生成。因而这就是关系“\(\alpha\) 与 \(\beta\) 位于偶 \((\alpha,\beta)\) 的框架的同一连通分支中”。

命题 5. —— 设 G' 为群胚，η, η' 为从 Coh(φ,ψ) 到 G' 的群胚态射，k 为从 Som(G) 到 Fl(G') 的映射。为使 k 是连接 η 与 η' 的同伦，必须且只需下列两个条件成立：

\begin{enumerate}
\def\labelenumi{(\roman{enumi})}
\item
  映射 \(k\) 是连接 \(\eta \circ \alpha\) 与 \(\eta' \circ \alpha\) 的同伦；
\item
  对 H 的每个顶点 a，有
\end{enumerate}

\[
\eta (h (a)) k (\psi (a)) = k (\varphi (a)) \eta^ {\prime} (h (a)).
\]

按定义，为使 k 是连接 \(\eta\) 与 \(\eta'\) 的同伦，必须且只需下列两个条件成立（回忆 \(\operatorname{Som}(G) = \operatorname{Som}(\operatorname{Coh}(\varphi, \psi))\) ）：

\begin{enumerate}
\def\labelenumi{\alph{enumi})}
\item
  对每个顶点 \(x\) （\(\operatorname{Coh}(\varphi, \psi)\) 的），\(k(x)\) 连接 \(\eta(x)\) 与 \(\eta'(x)\) ；
\item
  对每一对顶点 \((x,y)\) （属于 \(\mathrm{Coh}(\varphi,\psi)\) ）以及每条箭 \(f\in\mathrm{Fl}_{x,y}(\mathrm{Coh}(\varphi,\psi))\) ，有 \(\eta(f)k(y)=k(x)\eta'(f)\) 。
\end{enumerate}

由命题 2，只需验证条件 \(b\) ）——当 \(f\) 属于 \(\mathrm{Fl}(\alpha)\) 的像或 \(h\) 的像时；于是 \(b\) ) 等价于下列两个条件 \(c\) ) 与 \(d\) ) 的合取：

\begin{enumerate}
\def\labelenumi{\alph{enumi})}
\setcounter{enumi}{2}
\item
  当 \(x \in \operatorname{Som}(\mathrm{G})\) 、\(y \in \operatorname{Som}(\mathrm{G})\) 且 \(g \in \operatorname{Fl}_{x,y}(\mathrm{G})\) 时，有 \(\eta(\alpha(g))k(y) = k(x)\eta'(\alpha(g))\) ；
\item
  对每个 \(a \in \operatorname{Som}(\mathrm{H})\) ，有 \(\eta(h(a))k(\psi(a)) = k(\varphi(a))\eta'(h(a))\) 。条件 (i) 等价于 a) 与 c) 的合取，而条件 (ii) 即条件 d)，由此得推论。
\end{enumerate}

